% Folder 144, batch 05 (archivists' pages 81-100).
% Pass: Opus 5.5 (claude-opus-5-5), 2026-10-03 — first pass, unchecked against the pages by a human.
%
% All twenty pages are mathematics. His own pagination stands on the odd
% pages (17 on p. 81 to 26 on p. 99); the even pages are versos carrying the
% continuation. The batch opens in mid-argument (relation (59)) and closes on
% a sentence broken off (« Quand » on p. 100). Diagonal margin notes were read
% from rotated crops; many of their words remain marked.

\documentclass[11pt,a4paper]{article}
\input{../preamble/grothendieck}

\folder{144}
\batch{5}
\pages{81}{100}
\foldertitle{[Suite autour de Teichmüller dont] Teichmülleries, 1981-1982 : notes manuscrites (1981-1983, s.d.), lettres (1981, s.d.).}
\dating{1981-1983}
\watermark{Édition de démonstration}

\begin{document}

\page{81}
\note{p. 17 de l'auteur. La page commence au milieu d'un argument : la relation (59) et les notations (58) précèdent le lot.}

\[
  \text{(59)}\qquad \rho_{n-1}\rho_{n-2}\cdots\rho_1\rho_0 = 1 ,
\]
et on trouve en fait
\[
  \text{(60)}\qquad \Pi_P \simeq \Pi_{P_0} \simeq
  \{\rho_0,\rho_1,\ldots,\rho_{n-1} \mid \rho_{n-1}\rho_{n-2}\cdots\rho_1\rho_0 = 1\} ,
\]
où $P_0$ désigne le polygone combinatoire-type d'ordre $n$, i.e. muni d'un épinglage par un repère.

\uncertain{Venons-en} à $\Pi^\tau$, on introduit les
\[
  \text{(61)}\qquad \tau_a = (\lambda_a, \tau) \in \Pi_P^\tau \qquad (a \in A)
\]
on trouve que
\[
  \text{(62)}\qquad \Pi^\tau \simeq \{(\tau_a)_{a\in A} \mid \tau_a^2 = 1\} ,
\]
les relations entre les $\tau_a$ et les $\rho_s$ étant données par
\[
  \text{(63)}\qquad \rho_s = \tau_a \tau_{a'} ,
\]
où $a, a'$ sont les deux arêtes incidentes à $s$, \struck{\ill{}} avec \struck{\ill{}} $a = a(s,\omega)$ i.e. $a = \omega(a')$ i.e. $a' = \omega'(a)$.

Si on choisit \struck{\ill{}} une arête origine $a_0$ \add{($\to R_\infty$ …)} $a_0 = a(s_0,\omega)$, donc $(s_0,a_0)$ repère $\omega$-direct, \uncertain{et} \uncertain{pour}
\[
  \text{(64)}\qquad \tau_i = \tau_{\omega^i a_0} \qquad (i \in \mathbb{Z}/n\mathbb{Z})
\]
\marginal{Attention, le $\tau_{a_0}$ noté ici $\tau_0$, est noté $\tau_\infty$ dans (64) \ill{} \uncertain{des suites} \ill{} \uncertain{antérieures}, puis) en fait \ill{} (car $a_0$ est défini en \ill{} $R_\infty$ …)}
on trouve donc, avec les notations (58)
\[
  \text{(65)}\qquad \rho_i = \tau_i \tau_{i-1}
\]
(de sorte que la relation (59) est conséquence des relations $\tau_i^2 = 1$). L'écriture (62) devient
\[
  \text{(62 bis)}\qquad \Pi^\tau \simeq \{\tau_0,\ldots,\tau_{n-1} \mid \tau_0^2 = \tau_1^2 = \cdots = \tau_{n-1}^2 = 1\} .
\]

Considérons maintenant le stabilisateur de $P_\omega$ dans $D_P \times \{1,\tau\}$, il est formé des
\[
  \text{(66)}\qquad \tilde g^{0} = (g, \tau^{\alpha(g)}) \qquad
  \alpha(g) \in \mathbb{Z}/2\mathbb{Z},\quad
  \alpha(g) = \begin{cases} 0 & \text{si } \operatorname{sg}(g) = 1 \\ 1 & \text{si } \operatorname{sg}(g) = -1 \end{cases}
\]

\page{82}
On a alors un homomorphisme canonique, \struck{également} noté $g \mapsto \tilde g$ \struck{($g \mapsto \tilde g_\pi$ \ill{})}
\[
  \text{(67)}\qquad D \xrightarrow{\ \sim\ } (D \times \{1,\tau\})_{P_\omega} \longrightarrow \Pi_P^{D,\tau}
\]
qui peut s'expliciter comme
\[
  \text{(68)}\qquad g \longmapsto \tilde g = (1_{P_\omega}, \tilde g^{0})
\]
\note{un indice $\pi$ est biffé sous le $\tilde g$ de gauche.}
où $1_{P_\omega}$ est le chemin \struck{trivial} ponctuel de $P_\omega$ vers lui-même. On trouve en particulier
\[
  \text{(69)}\qquad \rho \overset{\text{déf}}{=} \tilde\omega_\pi \in \Pi_P^D \subset \Pi_P^{D,\tau} \qquad (\rho^n = 1)
\]
et des éléments
\[
  \text{(70)}\qquad \tilde\sigma_a = (\sigma_a^0)^{\sim} \in \Pi_P^{D,\tau} \quad \text{pour } a \in A \qquad (\tilde\sigma_a^2 = 1),
\]
où
\[
  \text{(71)}\qquad \sigma_a^0 \in D_P
\]
est l'unique \ill{} $\neq \mathrm{id}$ qui fixe $a$ (agissant sur $\Sigma$ comme la symétrie p.r. à l'axe de $\Sigma$ qui passe par $Q_a$).

\marginal{\uncertain{Je mets l'exposant} $\sim$, pour le distinguer de l'élément de $\Pi_P^D$ \ill{}, \uncertain{introduit ensuite} en (72)}

On pose alors
\[
  \text{(72)}\qquad \sigma_a = \tilde\sigma_a \tau_a = \tau_a \tilde\sigma_a = (\lambda_a, \sigma_a^0) \in \Pi_P^D \subset \Pi_P^{D,\tau}
\]
c'est donc un élément \emph{d'ordre 2} de $\Pi_P^D$.

\struck{On reconstitue les $\rho_s$} \add{Notons les formules}
\[
  \text{(73)}\qquad
  \begin{cases}
    \rho(\tau_a) = \tau_{\rho(a)} \\
    \rho(\tilde\sigma_a) = \tilde\sigma_{\rho(a)} \\
    \rho(\sigma_a) = \sigma_{\rho(a)}
  \end{cases}
\]
Introduisons aussi les
\[
  \text{(74)}\qquad \varepsilon_s = \sigma_a \rho = (\lambda_a, \sigma_s^0) \in \Pi_P^D
  \qquad s \in S, \text{ où } a = a(s,\omega)
\]
\note{après $\sigma_a$ un mot est biffé, et l'indice de $\lambda$ est corrigé en $a$.}
\[
  \text{(75)}\qquad \text{où } \sigma_s^0 \in D_P \text{ est l'unique \ill{} } \neq 1 \text{ de } P \text{ qui fixe le sommet } s \in S
\]
\marginal{NB On a $\varepsilon_s^{-1} = \rho^{-1}\sigma_a = (\lambda_{a'}, \sigma_s^0)$ (74), $a' = \omega' a$ l'autre arête incidente à $s$.}

\page{83}
\note{p. 18 de l'auteur.}
On a alors (cf. (56) pour la définition de $\rho_s$)
\[
  \text{(76)}\qquad \rho_s = \varepsilon_s^2 = \sigma_a \rho \sigma_a \rho
\]
\note{la suite de la ligne (76) est biffée et illisible.}
\marginal{NB $\rho_0 = \sigma_{a_0}\rho\sigma_{a_0}\rho$, \quad $\rho_1 = \rho(\rho_0) = \rho\sigma_{a_0}\rho\sigma_{a_0}$}
Cela \add{(76) et (73 c)} montre que $\rho, \sigma_{a_0}$ engendrent $\Pi_P^D$, et on trouve facilement
\[
  \text{(77)}\qquad \Pi_P^D \simeq (\rho, \sigma_{a_0} \mid \rho^n = \sigma_{a_0}^2 = 1) \simeq \mathbb{Z}/n\mathbb{Z} * \mathbb{Z}/2\mathbb{Z}
\]
(car la relation (59) est déjà conséquence, moyennant (76) et (73 c), des deux relations fondamentales
\[
  \text{(78)}\qquad \rho^n = 1, \quad \sigma_{a_0}^2 = 1 .
\]
\marginal{NB Pour $n=3$, on identifie \ill{} les notations \ill{} \ill{} \ill{} \ill{} $\sigma_{a_0}$}

Les formules (73) se complètent en
\[
  \text{(79)}\qquad
  \begin{cases}
    \rho(\varepsilon_s) = \varepsilon_{\rho(s)} \\
    \rho(\rho_s) = \rho_{\rho(s)}
  \end{cases}
  \qquad (s \in S)
\]
conséquence immédiate de (73), (74), (76), ce qui nous montre la façon dont $\rho$ opère sur $\Pi_P$.

\struck{Quant aux $\tilde\sigma_b$ $(b \in A)$, ils opèrent par transport de structure}
\[
  \text{\struck{(80)}}\qquad
  \text{\struck{$\tilde\sigma_b(\varepsilon_s) = \varepsilon_{\sigma_b^0(s)}^{-1}$}}, \qquad
  \text{\struck{$\tilde\sigma_b(\rho_s) = \rho_{\sigma_b^0(s)}^{-1}$}}, \qquad
  \text{\struck{$\forall s \in S,\ b \in A$}}
\]
\struck{\ill{} trouvons l'opération}
\note{passage encadré et barré en croix ; des flèches le relient aux formules qui suivent.}

\marginal{À rectifier : $\tau_a(\rho_s) = \rho_s^{-1}$ si $s$ incident à $a \in A$}

On trouve \ill{} (\uncertain{directement} sur (65) et (62 bis))
\[
  \text{(80)}\qquad
  \begin{aligned}
    &\tau_{a_0}(\rho_0) = \rho_0^{-1}, \quad \tau_{a_0}(\rho_1) = \rho_1^{-1}, \quad \tau_{a_0}(\rho_2) = \rho_1^{-1}(\rho_2^{-1}) \quad \cdots \\
    &\tau_{a_0}(\rho_i) = (\rho_{i-1}\rho_{i-2}\cdots\rho_1)^{-1}(\rho_i^{-1}) \\
    &\phantom{\tau_{a_0}(\rho_i)} = (\rho_{n-1}\cdots\rho_{i+1})(\rho_i^{-1}) \qquad 0 \leq i \leq n-1
  \end{aligned}
\]
\note{sous le dernier produit, une remarque à peine lisible : « \ill{} avec $\rho_0$, pas $\rho_{n-1}$ ! ».}
\marginal{Attention, j'ai pris ici $\tau_{a_0}$ au lieu de $\tau_0$, (80) car pour $n=3$ la « notation reçue » pour $\tau_{a_0}$ est $\tau_\infty$, et non $\tau_0$ !}
\marginal{NB L'action de $\tau_a$ sur $\Pi_P$ est la composée de $\tau : \pi_1(\Sigma^*, P_\omega) \to \pi_1(\Sigma^*, P_{\omega'})$ avec $\lambda_a : \pi_1(\Sigma^*, P_{\omega'}) \to \pi_1(\Sigma^*, P_\omega)$}

\struck{Pour} Les relations (79) se généralisent ainsi : en
\[
  \text{(81)}\qquad
  \begin{cases}
    \tilde g(\varepsilon_s) = (\varepsilon_{g^0(s)})^{\operatorname{sg}(g)} \\
    \tilde g(\rho_s) = (\rho_{g^0(s)})^{\operatorname{sg}(g)}
  \end{cases}
  \qquad (g \in D_P,\ s \in S),
\]
\marginal{NB L'action de $\tilde D_P \subset \Pi_P^{D,\tau}$ sur $\Pi$ est par « transport de structure » via celle de $D_P^0 \subset D_P \times \{1,\tau\}$ sur $(\Sigma^*, P_\omega)$}

\page{84}
(et en particulier, si $g$ de la forme $\sigma_a^0$, on trouve
\[
  \text{(82)}\qquad
  \begin{cases}
    \tilde\sigma_a(\varepsilon_s) = \varepsilon_{\sigma_a^0(s)}^{-1} \\
    \tilde\sigma_a(\rho_s) = \rho_{\sigma_a^0(s)}^{-1}
  \end{cases}
  \qquad s \in S,\ a \in A
\]
ce qui donne l'opération des $\tilde\sigma_a$ sur $\Pi_P$, \struck{\ill{}} d'où \add{en principe} \struck{\ill{}} (avec celle des $\tau_a$) l'opération des \ill{} sur $\Pi_P$ \add{(comme $\sigma_{a_0} = \tilde\sigma_{a_0}\tau_{a_0}$)}. En particulier,
\[
  \begin{aligned}
    &\sigma_{a_0}(\rho_0) = \rho_1, \quad \sigma_{a_0}(\rho_1) = \rho_0, \quad
    \sigma_{a_0}(\rho_i) = \tau_{a_0}(\rho_{n+1-i}^{-1}) \\
    &\phantom{\sigma_{a_0}(\rho_i)} = (\rho_{i-1}\cdots\rho_{n+3-i})(\rho_{n+1-i}) \qquad \text{pour } 2 \leq i \leq n-1 .
  \end{aligned}
\]
\note{les indices de cette dernière ligne sont serrés et de lecture douteuse.}
\marginal{d'où $\tilde\sigma_{a_0}(\rho_i) = \rho_{1-i}^{-1}$}
\marginal{NB L'action de $\sigma_a$ sur $\Pi_P$ est la composée de $\pi_1(\sigma_a^0) : \pi_1(\Sigma^*, P_\omega) \to \pi_1(\Sigma^*, P_{\omega'})$ avec $\lambda_a : \pi_1(\Sigma^*, P_{\omega'}) \to \pi_1(\Sigma^*, P_\omega)$ \ill{}}

\struck{Or} On trouve d'ailleurs \struck{$\Pi_P^\tau$}
\[
  \text{(84)}\qquad \Pi_{P_\omega}^{D,\tau} \simeq \Pi_P^D \cdot \{1, \tau_{a_0}\}
\]
produit \uncertain{semi-direct}, il faut expliciter encore l'opération de $\tau_{a_0}$ sur $\Pi_P^D$ i.e. sur $\rho$ et sur $\sigma_{a_0}$, on a
\[
  \text{(85)}\qquad \tau_{a_0}(\rho) = \rho^{-1}, \quad \tau_{a_0}(\sigma_{a_0}) = \sigma_{a_0} ,
\]
\note{dans (85), un $\sigma_{a_0}$ est biffé devant $(\rho^{-1})$ ; il n'y a pas de (83) sur la page.}
la deuxième relation est déjà notée (72), la première s'obtient en écrivant
\[
  \tau_{a_0}\underbrace{\rho\,\tau_{a_0}\rho^{-1}}_{\tau_{a_1}} = \tau_{a_0}\tau_{a_1} = \rho_1^{-1}
  = (\rho\sigma_{a_0}\rho\sigma_{a_0})^{-1} = \sigma_{a_0}\rho^{-1}\sigma_{a_0}\rho^{-1}
\]
\note{l'accolade tombe sous $\rho\tau_{a_0}\rho^{-1}$ ; un $\rho$ initial est biffé.}
d'où (85) ; \struck{qui s'écrivent} \add{plus simplement} elles s'écrivent
\[
  \tilde\sigma_{a_0}(\rho) = \rho^{-1}
\]
qui est trivial, puisque \struck{pour} $\rho = \tilde\omega$ \struck{\ill{}} et que $g \mapsto \tilde g$ est un homomorphisme de groupes $D_P \to \Pi_1^{D,\tau}$, enfin $\sigma_{a_0}^0(\rho) = \rho^{-1}$ …. Je \ill{} plus \ill{} peut-être d'erreur
\[
  \text{(86)}\qquad
  \begin{cases}
    \Pi_P^{D,\tau} \simeq \Pi_P^D \cdot \{1, \tilde\sigma_{a_0}\} \\
    \tilde\sigma_{a_0}(\rho) = \rho^{-1}, \quad \tilde\sigma_{a_0}(\sigma_{a_0}) = \sigma_{a_0}
  \end{cases}
\]
\[
  \Bigl[ = \{\rho, \sigma_{a_0}, \tilde\sigma_{a_0} \mid \rho^n = \sigma_{a_0}^2 = \tilde\sigma_{a_0}^2 = (\sigma_{a_0}\tilde\sigma_{a_0})^2 = \tilde\sigma_{a_0}(\rho)\rho = 1\} \Bigr]
\]


\page{85}
\note{p. 19 de l'auteur.}
ce qui s'écrit aussi
\[
  \text{(87)}\qquad \Pi_P^{D,\tau} \simeq D(\rho) \underset{D_1}{*} D(\sigma_a)
\]
somme amalgamée de groupes, sous $D_1 = \{1, \tilde\sigma_a\}$, où
\[
  \text{(88)}\qquad
  \begin{cases}
    D(\rho) = D^+ \cdot \{1, \tilde\sigma_a\} \\
    D(\sigma_a) = \{1, \sigma_a\} \times \{1, \tilde\sigma_a\} \\
    D_1 = \{1, \tilde\sigma_a\}
  \end{cases}
\]
avec, pour $D(\rho)$ : groupe diédral d'indice $n$ \struck{engendré} \add{défini} par $\rho$ (sous $D^+$ : « engendré par $\rho$, avec $\rho^n = 1$ ») ; pour $D(\sigma_a)$ : groupe diédral d'indice 2 défini par $\sigma_a$.
\note{à la première ligne de (88), un $D(\tilde\sigma_a)$ est biffé après $\{1,\tilde\sigma_a\}$.}

L'écriture (77) s'écrit aussi
\[
  \text{(89)}\qquad \Pi_P^D \xrightarrow{\ \sim\ } \Pi_{0,3}/\langle \rho_1^2, \rho_\infty^n \rangle
  = \pi_1\bigl(\mathbb{P}^1(\mathbb{C}), \underbrace{\infty\cdot(0) + 2(1) + n(\infty)}_{\text{signature}} ; P_+\bigr)
\]
en transformant $(\rho, \sigma_{a_0})$ en $(\rho_\infty, \rho_1)$. Une façon a priori d'obtenir un isomorphisme (89) est de considérer la courbe complexe quotient de $\Sigma_P$ par l'action de $D_P$, \struck{qui} action \add{fidèle} qui a comme points de ramification les $R_s$ (ramification 2), les $Q_s$ (ramification 2), et $P_\omega, P_{\omega'}$ (ramification $n$) — le quotient est \add{une courbe alg. complexe} isomorphe à $\mathbb{P}^1_{\mathbb{C}}$, l'isomorphisme étant uniquement déterminé quand on le normalise par
\[
  \text{(90)}\qquad
  \begin{cases}
    R_s \longmapsto 0 & (s \in S) \\
    Q_a \longmapsto 1 & (a \in A) \\
    P_\omega \longmapsto \infty & (\omega \in \underline{\omega})
  \end{cases}
  \qquad (\Sigma_P \to \mathbb{P}^1_{\mathbb{C}}) .
\]
On peut \uncertain{aussi} \add{préciser} que la carte envisagée sur $\Sigma_P$ est l'image inverse de la carte universelle, par un morphisme $\Sigma_P \to \mathbb{P}^1_{\mathbb{C}}$ bien déterminé, \struck{qui} [Dans le cas épinglé
\marginal{C'est un peu une tricherie, car pour construire cette carte, \ill{} l'explicitation qui \ill{} \ill{} point \ill{}, il \uncertain{conviendrait} \ill{} \ill{} morphisme $\Sigma_P \to \mathbb{P}^1_{\mathbb{C}}$, cf. plus bas}

\page{86}
\[
  \text{(91)}\qquad
  \left\{
  \begin{array}{l}
    P = P_0 = \text{polygone comb. formé des } S_0 = \mu_n(\mathbb{C}) \\
    \text{et des arcs de cercle sur } U = \{z \in \mathbb{C} \mid |z| = 1\} \\
    \text{qui les joignent (car } P \text{ \emph{épinglé})}
  \end{array}
  \right.
\]
\[
  \Sigma_P = \mathbb{P}^1_{\mathbb{C}} = \hat{\mathbb{C}}, \qquad
  D_P^+ = \mu_n(\mathbb{C}) \text{ opérant par homothéties}, \qquad
  \tau(z) = 1/\bar z
\]
\struck{L'application} On peut considérer que la carte envisagée \add{$X_n$} est l'image inverse, par
\[
  z \longmapsto z^n \qquad \mathbb{P}^1(\mathbb{C}) \to \mathbb{P}^1(\mathbb{C})
\]
de la carte cyclotomique d'ordre 1, $X_1$, déjà considérée (p. 7), ayant un seul sommet, un seul arc et deux faces (qui sont des monogones) (\uncertain{cercle} \ill{} $-1$, qui sépare les faces $X_1^+$ et $X_1^-$ contenant $0$ et $\infty$), qui est image inverse de la carte universelle $X_0 : \{0\} \subset [0,1] \subset \mathbb{P}^1_{\mathbb{C}}$, par l'application
\[
  z \longmapsto -\frac{1}{4z}(z-1)^2 \qquad \mathbb{P}^1(\mathbb{C}) \to \mathbb{P}^1(\mathbb{C})
\]
\drawing{petite sphère en marge : $0$ au pôle supérieur, $\infty$ au pôle inférieur, l'équateur marqué des points $1$ et $-1$.}
(cf. (23)) de sorte que l'application cherchée $X_n \to X_0$ est
\[
  \text{(92)}\qquad z \longmapsto -\frac{1}{4z^n}(z^n - 1)^2 \qquad
  \underbrace{X_n}_{\substack{\text{carte cyclotomique}\\ \text{type d'ordre } n}} \xrightarrow{\ g_n\ }
  \underbrace{X_0}_{\substack{\text{carte}\\ \text{universelle}}} . \;]
\]
\note{le crochet fermant clôt le passage ouvert p. 85 par « [Dans le cas épinglé ». Au-dessus de $g_n$, un mot est biffé.}

Ainsi, $X_n$ devient un revêtement galoisien de $X_0$, de groupe de Galois $\mathbb{D}_n = D_{P_0}$, ramifié comme il se doit sur les seuls points $0, 1, \infty$, avec en ces points la ramification $2, 2, n$ respectivement, et $\Sigma_{P_0}^* = X_n^*$ devient un revêtement fini ramifié \add{galoisien} de $X_0 \setminus \{0\}$, ramifié sur les seuls points $1, \infty$, avec ramification $2, n$ en ces

\page{87}
\note{p. 20 de l'auteur.}
points. D'après la théorie générale des groupes fondamentaux mixtes et de leurs relations avec les groupes fondamentaux à signatures, on en tire donc un isomorphisme (extérieur pour le moment) canonique
\[
  \text{(93)}\qquad
  \begin{aligned}
    \pi_1(X_n^*, P_{n+}) &\xrightarrow{\ \sim\ } \pi_1\bigl(\underbrace{X_0^*}_{X_0 \setminus \{0\}}, 2\{1\} + n\{\infty\} ; \underbrace{P_{0+}}_{\exp 2i\pi/6}\bigr) \\
    &\simeq \pi_1\bigl(\underbrace{X_0}_{\mathbb{P}^1(\mathbb{C})}, \infty\{0\} + 2\{1\} + n\{\infty\} ; P_{0+}\bigr)
  \end{aligned}
\]
\note{sous $P_{n+}$ est écrit « $= 0$ » (lecture douteuse) ; après $X_0$ un exposant est surchargé.}
Pour trouver un isomorphisme pas extérieur, il faut changer le point base $P_{n+} = \ill{}$ en un autre, qui s'envoie dans \struck{\ill{}} $X_0^* \setminus \{1,\infty\} = X_0 \setminus \{0,1,\infty\} = U_{0,3}$,

\drawing{deux sphères reliées par une flèche. À gauche, $X_n$ : pôles $P_\omega$ (haut) et $P_{\omega'}$ (bas), sur l'équateur $R_0$, $Q_{a_0}$, $R_1$ ; méridiens ; le point $P$ près de $P_\omega$, avec les chemins $\alpha$, $\alpha_\rho$ vers $\rho(P)$, $\rho(\alpha)$, et sous l'équateur $\alpha_\sigma$, $\sigma_{a_0}^0(P)$, $\sigma_{a_0}^0(\alpha)$, $\lambda_{a_0}$. À droite, $X_0$ : le pôle $P_+$, le point $\infty$, le segment $[0,1]$ sur l'équateur, les hémisphères $X_0^+$ et $X_0^-$, et le lacet $g_n(\alpha_\rho)$.}

On considère la subdivision barycentrique de la carte $X_n$, image inverse de la subdivision barycentrique de $X_0$, dont les deux faces sont les deux hémisphères $X_0^+$ et $X_0^-$ — les images inverses de celles-ci sont les $4n$ « repères » \add{ou « drapeaux »} de la subdivision barycentrique de $X_n$, dont les $2n$ repères \emph{directs} \struck{qui} correspondant aux $2n$ repères du polygone $P$. L'épinglage de ce dernier nous fournit un repère origine
\[
  \text{(94)}\qquad r_0 = (R_0, Q_{a_0}, P_\omega) \qquad \bigl(\xrightarrow{\ g_n\ } X_0^+\bigr)
\]
qui par $g_n$ est appliqué homéomorphiquement sur $X_0^+$, avec \struck{$R_0$}
\[
  g_n : \quad R_0 \longmapsto 0, \quad Q_{a_0} \longmapsto 1, \quad P_\omega \longmapsto \infty
\]
Prenons le point $P \in r_0$ qui s'envoie sur le « pôle nord » $P_+$ $(= \exp 2i\pi/6) \in X_0^+$,
\[
  g_n : \quad P \longmapsto P_+ .
\]

\page{88}
La théorie générale donne donc un iso (pas extérieur !)
\[
  \text{(95)}\qquad
  \begin{aligned}
    \pi_1(X_n^*, \text{\add{$\mathbb{D}_n;$}}\, P) &\xrightarrow{\ \sim\ } \pi_1\bigl(X_0^*, 2\{1\} + n\{\infty\} ; P_+\bigr) \qquad X_0^* \overset{\text{déf}}{=} X_0 \setminus \{0\} \\
    &\simeq \Pi_{0,3}/\langle \rho_1^2, \rho_\infty^n \rangle ,
  \end{aligned}
\]
pour obtenir un isomorphisme précis (93), il faut composer (95) avec un isomorphisme
\[
  \text{(96)}\qquad \pi_1(X_n^*, \text{\add{$\mathbb{D}_n;$}}\, P_+) \xrightarrow{\ \sim\ } \pi_1(X_n^*, \mathbb{D}_n ; P)
\]
déduit d'une classe de chemins $P_+ \to P$ dans $X_n^*$ — on prendra la \add{unique} classe de chemins provenant d'un chemin $\alpha$ \emph{dans $r_0$}, comme choix le plus naturel, ce qui précise 96. L'homomorphisme (95), quand à lui, est donné par la formule des plus évidentes
\[
  (l, g) \longmapsto g_n(l)
\]
où $g \in \mathbb{D}_n$, $l : g(P_*) \to P_*$ [chemin dans $X_n^*$, ou mieux, chemin dans $X_n^{**} = X_n^* \setminus \{(Q_a), P_\omega, P_{\omega'}\} = X_n \setminus \{\text{pts de ramification}\} = X_n \mid U_{0,3}$], d'où $g_n(l) : g_n(P) \to g_n(P)$ un lacet dans $P_+ = g_n(P)$ dont on prend la classe dans le quotient de $\Pi_{0,3}$ — qui ne dépend pas des choix faits, mais seulement de l'élément de $\pi_1(X_n^*, \mathbb{D}_n ; P)$ représenté par $(l, g)$ ….
\note{l'indice de $P_*$ est surchargé ; on attendrait $P$.}

Considérons donc l'application \add{nous intéresse \ill{}} composée des (95) (96), qui
\[
  \text{(97)}\qquad
  \begin{aligned}
    \pi_1(X_n^*, \mathbb{D}_n ; P_+) &\xrightarrow{\ g_n\ } \pi_1\bigl(X_0, \infty\cdot\{0\} + 2\cdot\{1\} + n\{\infty\} ; P_+\bigr) \\
    &\simeq \Pi_{0,3}/\langle \rho_1^2, \rho_\infty^n \rangle
  \end{aligned}
\]
et calculons ses valeurs pour $\rho, \sigma_{a_0}$. L'image de $\rho$ \struck{par} (96), est donnée par le chemin
\[
  \rho(P) \xrightarrow{\ \alpha_\rho\ } P \qquad \text{dans } X_n^{**}
\]
qui franchit le 1-squelette de la subdivision
\marginal{NB $\alpha_\rho$ est obtenu par déformation dans $X_n^*$ du chemin composé $\alpha \circ \rho(\alpha)^{-1}$}

\page{89}
\note{p. 21 de l'auteur. Les arcs $(\infty,0)$, $(\infty,1)$, $(0,1)$, $(1,\infty)$ sont surmontés d'un arc sur la page.}
barycentrique juste deux fois, \struck{\ill{}} en traversant respectivement les arêtes $P_\omega R_1$ $\bigl(\xmapsto{g_n} (\infty,0)\bigr)$ et $P_\omega Q_{a_0}$ $\bigl(\xmapsto{g_n} (\infty,1)\bigr)$ — donc $g_n(\alpha_\rho)$ est un lacet dans $U_{0,3} = X_0^* \setminus \{1,\infty\}$ \struck{qui} en $P_+$, qui \struck{franchit} traverse une fois l'arc $(\infty,0)$ de l'équateur, une fois l'arc $(\infty,1)$, dans cet ordre. Sa classe dans $\Pi_{0,3}$ est donc $\rho_\infty^{-1}$ :
\[
  (*)\qquad g_n(\rho) = \rho_\infty^{-1} .
\]
\struck{L'élément de} \add{L'image de} $\sigma_{a_0} = (\lambda_{a_0}, \sigma_{a_0}^0)$ par (96) est représentée par $\sigma_{a_0}^0(P) \to P$, qui est le composé dans $\Pi_1(X_n^*)$ \struck{\ill{}} $\alpha\, \lambda_{a_0} \sigma_{a_0}^0(\alpha)^{-1}$, lequel chemin peut se déformer en le chemin \struck{\ill{}}
\[
  \alpha_\sigma : \sigma_{a_0}^0(P) \longrightarrow P
\]
situé tout entier dans $X_n^{**}$, qui franchit le 1-squelette \add{de la} subdivision barycentrique en exactement deux pts, qui sont respectivement sur les arcs $P_\omega Q_{a_0}$ $\bigl(\xmapsto{g_n} (\infty,1)\bigr)$ et $(Q_{a_0}, R_0)$ $\bigl(\xmapsto{g_n} (0,1)\bigr)$, donc $g_n(\alpha_\sigma)$ est un lacet en $P_+$ qui franchit $(\infty,1)$ puis $(0,1)$ dans cet ordre, donc qui \struck{est} a la classe $\rho_1^{-1}$. \struck{On} On aurait pu prendre aussi $\alpha'_\sigma : \sigma_{a_0}^0(P) \to P$, franchissant d'abord $(Q_{a_0}, R_1)$ $\bigl(\xmapsto{g_n} (0,1)\bigr)$ puis $(Q_{a_0}, P_\omega)$ $\bigl(\xmapsto{g_n} (1,\infty)\bigr)$ — \struck{on aurait} trouverait $g_n(\alpha'_\sigma) = \rho_1$ — mais dans le quotient $\Pi_{0,3}/\langle \rho_1^2, \rho_\infty^n \rangle$, $\rho_1$ et $\rho_1^{-1}$ sont égaux, comme il se doit. Donc on trouve \uncertain{aussi},
\[
  (**)\qquad g_n(\sigma_{a_0}) = \rho_1 \quad (= \rho_1^{-1})
\]

\page{90}
et en comparant (*) et (**) on trouve
\[
  g_n(\underbrace{\varepsilon_0}_{\sigma_{a_0}\rho}) = g_n(\sigma_{a_0})\, g_n(\rho) = \rho_1^{-1}\rho_\infty^{-1} = (\rho_\infty\rho_1)^{-1} = \rho_0
\]
En résumé, on trouve donc
\[
  \text{(98)}\qquad
  \begin{cases}
    g_n : \pi_1(X_n^*, \mathbb{D}_n ; P_{n+}) \xrightarrow{\ \sim\ } \Pi_{0,3}/\langle \rho_1^2, \rho_\infty^n \rangle \\
    \qquad\qquad = \pi_1\bigl(\underbrace{X_0}_{\mathbb{P}^1_{\mathbb{C}}}, \underbrace{\infty\{0\} + 2\{1\} + n\{\infty\}}_{\text{signature } \nu_n} ; P_+\bigr) \\
    \varepsilon_0 \longmapsto \rho_0 \\
    \sigma_{a_0} \longmapsto \rho_1 = \rho_1^{-1} \\
    \rho \longmapsto \rho_\infty^{-1}
  \end{cases}
\]

Par le même principe, on trouve un isomorphisme
\[
  \text{(99)}\qquad
  \begin{aligned}
    g_n : \pi_1(X_n^*, \mathbb{D}_n \times \{1,\tau\} ; P_{n+}) &\xrightarrow{\ \sim\ } \pi_1(X_0, \nu_n, \{1,\tau\} ; P_+) \\
    \simeq \Pi_n^{D,\tau} \simeq \Pi_n^D \cdot \{1, \tilde\sigma_{a_0}\}
    &\simeq \bigl(\Pi_{0,3}/\langle \rho_1^2, \rho_\infty^n \rangle\bigr) \cdot \{1, \ill{}\}
  \end{aligned}
\]
\note{dans (99), un symbole est biffé avant $\nu_n$ ; le dernier facteur est surchargé.}
\marginal{NB $\tau_\infty(\rho_0) = \rho_0^{-1}$, $\tau_\infty(\rho_1) = \rho_1^{-1}$, $\tau_\infty(\rho_\infty) = \rho_0(\rho_\infty^{-1})$ cf. (80) dans le cas $n=3$, ici $\tau_\infty = \tau_{a_0}$ par définition …}
qui prolonge (97) = (98), et qui se calcule de la même façon. Il suffit, pour achever de déterminer $g_n$, de calculer sa valeur sur \struck{\ill{}}
$\tilde\sigma_{a_0} = (1_{P_+}, \tilde\sigma_{a_0}^0)$. Ici $\tilde\sigma_{a_0}^0$ agit sur $X_n$ par la symétrie par rapport au grand cercle $P_\omega Q_{a_0} P_{\omega'}$, \struck{donc} son image dans $\pi_1(X_n^*, \mathbb{D}_n \times \{1,\tau\} ; P)$ est définie par le chemin \struck{\ill{}}
\[
  P' = \tilde\sigma_{a_0}^0(P) \longrightarrow P
\]
donné par $\alpha \circ \alpha'^{-1}$, où $\alpha' = \tilde\sigma_{a_0}^0(\alpha)$, ce chemin est homotope dans $X_n^*$ au chemin $\tilde\alpha_\sigma$, \struck{qui} dans $X_n^{**}$, qui franchit la subdivision barycentrique une seule fois, en $P_+ Q_{a_0}$ $\bigl(\xmapsto{g_n} (\infty,1)\bigr)$.
\drawing{en marge, une sphère : le pôle $P_+$, sur l'équateur $R_0$, $Q_{a_0}$, $R_1$ ; les chemins $\alpha$ et $\tilde\sigma_{a_0}^0(\alpha) = \alpha'$ descendent vers $P$ et $P' = \tilde\sigma_{a_0}^0(P)$, reliés par $\tilde\alpha_\sigma$.}

\page{91}
\note{p. 22 de l'auteur.}
Par $g_n$ ce chemin devient un chemin dans $U_{0,3}$
\[
  g_n(\tilde\alpha_\sigma) : \quad g_n(P') = g_n(\tilde\sigma_{a_0}^0(P)) = \tau\, g_n(P) = \tau(P_+) = P_- \longrightarrow g_n(P) = P_+
\]
qui franchit l'équateur \struck{\ill{}} exactement une fois, en traversant l'arc $(1,\infty)$ — c'est donc $\tau_{a_1}$, noté aussi $\tau_0$ (car $a_1$ est opposé au sommet $0$) donc on trouve
\[
  (**)\qquad g_n(\tilde\sigma_{a_0}) = \tau_0 \in \Pi_{0,3}^\tau/\langle \rho_1^2, \rho_\infty^n \rangle
\]
\drawing{petite sphère en marge : sur l'équateur $0$, $a_0$, $1$ ; derrière, $a_\infty$, $\infty$, $a_1$.}
Pour se rassurer, on va vérifier que conformément à (86), $\tilde\sigma_{a_0}(\rho) = \rho^{-1}$, $\tilde\sigma_{a_0}(\sigma_{a_0}) = \sigma_{a_0}$, on a bien
\[
  \tau_0(\rho_\infty^{-1}) = \rho_\infty, \qquad \tau_0(\rho_1) = \rho_1^{-1}
\]
\note{l'exposant de $\rho_\infty$ dans le second membre est surchargé.}
ce qui est contenu en effet dans les \ill{} (deux premières) formules (80) pour $n=3$, en prenant ici \struck{\ill{}} $a_0$ le côté $(1,\infty)$ du triangle $(0,1,\infty)$ dans $X_0$ …. On tire de \struck{\ill{}} (*)
\[
  g_n(\tau_{a_0}) = g_n(\tilde\sigma_{a_0} \cdot \sigma_{a_0}) = g_n(\tilde\sigma_{a_0})\, g_n(\sigma_{a_0}) = \tau_0 \rho_1 = \tau_\infty
\]
puisque $\rho_1 = \tau_0 \tau_\infty$, donc en résumé
\[
  \text{(100)}\qquad g_n(\tilde\sigma_{a_0}) = \tau_0, \qquad g_n(\tau_{a_0}) = \tau_\infty .
\]

\textbf{7.} Dans toutes ces considérations \add{du §6}, on a un peu perdu de vue les cartes, dont il n'a plus été question. Pourtant, le groupe $\Pi_P^\tau$ s'interprète comme un des groupes « cartographiques pondéré \struck{\ill{}}-polygonal », qui est tel que


\page{92}
\note{en tête de page, la marque « B ».}
les \struck{\ill{}} ensembles finis où opère $\Pi_P^\tau$ — i.e. munis d'endomorphismes \struck{\ill{}} involutifs $\tau_a$, indexés par les \emph{arêtes} de $P$ — s'interprètent en termes de telles cartes, mais en supposant que les $\tau_a$ opèrent sans points fixes. \struck{\ill{}} Si on associe à la carte $X$ $P$-pondérée l'ensemble \add{$R(X)$} de ses \emph{faces}, on trouve en effet un ensemble où des opérations $\tau_a$ \add{où $\tau_a$} correspondant à la symétrie p.r. \struck{\ill{}} : l'arête de type $a$ d'une face. Ici, la pondération (disons, \add{\emph{orientée} d'une} numérotation circulaire) des \struck{\ill{}} \add{faces} pourrait plutôt commencer par celle des sommets. Le foncteur
\marginal{Une pondération de type $P$}
\[
  \text{(101)}\qquad
  \begin{array}{ccc}
    X & \longmapsto & R(X) \\[2pt]
    \begin{array}{c}\text{cartes $P$-pondérées}\\ \text{finies}\\ \text{(cat. isotopique)}\end{array}
    & \xrightarrow{\ \approx\ } &
    \begin{array}{c}\Pi_P^\tau\text{-ensembles finis}\\ \text{spéciaux}\\ \text{(i.e. où les $\tau_i$}\\ \text{opèrent ss pts fixes)}\end{array}
  \end{array}
\]
\drawing{en marge gauche, un assemblage de faces quadrangulaires numérotées (côtés marqués $0$, $1$, $2$, $3$) autour d'une face centrale $r$ munie d'une flèche de rotation ; les faces voisines portent $\tau_0 r$, $\tau_1 r$, $\tau_2 r$, $\tau_3 r$.}
\note{l'étiquette (101) est surchargée sous le dessin.}
est une équivalence de catégories. \add{Elle} \struck{L'équivalence} \add{induit une équivalence}
\struck{\ill{} \ill{} \ill{} \ill{}}
\[
  \text{(102)}\qquad
  \begin{array}{ccc}
    \begin{array}{c}\text{cartes $P$-pondérées}\\ \text{finies \emph{orientées}}\end{array}
    & \xrightarrow{\ \approx\ } & \Pi_P\text{-ens. finis} \\[2pt]
    X & \longmapsto & R^+(X)
  \end{array}
\]
\note{devant (102), un $\Pi_P^\tau$ suivi d'un mot est biffé, ainsi que « \struck{statiques} » après « $P$-pondérées ».}
en associant à chaque telle carte l'ensemble $R^+(X)$ de ses \struck{repères} faces « directes », i.e. \struck{\ill{}} l'orientation induite par $X$ et celle définie par $\omega \in \underline{\omega}$, définissant un ordre circulaire sur l'une des sommets (ou des arêtes)

\page{93}
\note{p. 23 de l'auteur.}
de la face. On notera que pour décrire le groupe $\Pi_P^\tau$, et l'équivalence (101), on n'a pas besoin d'autre donnée que du pol. combinatoire $P$, le groupe $\Pi_P^\tau$ est défini simplement par (62) (sans référence au choix d'un \struck{point}-base $P_\omega$ dans un espace $\Sigma_P^*$ !). Par contre, pour définir $\Pi_P$ et l'équivalence (102), on a besoin du choix d'une orientation \uncertain{ω} de $P$, ce qui permet d'introduire $\Pi_P$ p.ex. comme $\pi_1(\Sigma_P^*, P_\omega)$. À vrai dire, on pourrait définir $\Pi_P$, sans choix de $\omega$, comme noyau de l'homom. can. $\Pi_P^\tau \to \{\pm 1\}$, de façon à avoir la suite exacte
\[
  \text{(103)}\qquad
  \begin{array}{ccccccccc}
    1 & \to & \Pi_P & \to & \Pi_P^\tau & \longrightarrow & \{\pm 1\} & \longrightarrow & 1 \\
      &     &       &     & \tau_a & \longmapsto & -1 & &
  \end{array}
\]
mais il reste que le choix d'un foncteur (102) dépend du choix d'une orientation $\omega$ de $P$.

Un foncteur quasi-inverse de (102) s'obtient en interprétant le deuxième membre par la théorie de Galois des revêtements étales de $\Sigma_P^*$
\[
  \text{(104)}\qquad
  \Pi_P\text{-ens. finis} \xrightarrow{\ \approx\ }
  \begin{array}{l}\text{revêtements ramifiés}\\ \text{finis de } \Sigma_P \text{, étales}\\ \text{au-dessus de } \Sigma_P^* ,\end{array}
\]
si on a un tel revêtement \struck{\ill{}} $X$ de $\Sigma_P$, l'image inverse de la carte \struck{\ill{}} \add{canonique} sur $\Sigma_P$ (jouant le rôle de \emph{carte $P$-pondérée orientée universelle})

\page{94}
\note{en tête de page, un chiffre biffé.}
est la carte $P$-pondérée \struck{cherchée} \add{orientée} cherchée.

Pour avoir un foncteur quasi-inverse de (101), on procède de façon analogue, en utilisant une orientation de $P$. Ce \uncertain{faisant}, on utilise l'un $\omega$ pour \uncertain{tordre} le corps $\mathbb{C}$, et le remplacer par $\mathbb{C}(\omega)$, et on considère \struck{la sphère} $\mathbb{C}(\omega)$-courbe $\Sigma_P(\mathbb{C}(\omega))$ au lieu de $\Sigma_P(\mathbb{C})$, qui \uncertain{a} \uncertain{aussi} un point canonique $P_+$ pouvant servir de pt base). On \struck{passe} interprète $\Pi_P^\tau$ comme $\pi_1(\Sigma_P^*, \{1,\tau\} ; P_\omega)$, et on procède comme p. 9' (dans le cas $n=3$, dont \ill{} ici une généralisation …).

Sûrement on doit pouvoir généraliser aussi \ill{} quelconque, les considérations du §5 (p. 12'-14), qui précisent (en relation avec \struck{\ill{}} les iso. (97) (99)) les relations entre cartes $P$-pondérées, et cartes $n$-gonales (i.e. dont toutes les faces sont $n$-gonales …). Je ne vais pas détailler ce point ici.

\page{95}
\note{p. 24 de l'auteur. Ici commence un nouveau paragraphe ; sa numérotation des formules repart de (1).}
\textbf{8.} Une façon d'obtenir une configuration « polygone régulier à $n$ côtés » \add{$P_n$} dans une sphère complexe \add{$\Sigma$}, \struck{\ill{}} i.e. dans une droite projective complexe, est de partir d'un polygone régulier à $2n$ côtés $\tilde P_{2n}$ (combinatoire, un revêtement de degré 2 de $P_n$) plongé dans un \struck{espace} \add{plan} vectoriel réel $V_{\tilde P_{2n}}$ (centre du polygone régulier) et de poser
\marginal{$P_n = P = (S, A, R \subset S \times A)$}
\[
  \text{(1)}\qquad \Sigma = \mathbb{P}(V_{\tilde P_{2n}})(\mathbb{C})
\]
en prenant
\[
  \text{(2)}\qquad P_n = \tilde P_{2n}/\{\pm 1\}
\]
\marginal{NB $-1$ opère sur le polygone combinatoire $\tilde P_{2n}$ par antipodisme}
\struck{et on} donc
\[
  S = \tilde S/\pm 1, \qquad A = \tilde A/\pm 1, \qquad R = \tilde R/\pm 1 ,
\]
et on définit le plongement
\[
  \text{(3)}\qquad S \hookrightarrow \Sigma
\]
par passage au quotient en termes de
\[
  \begin{array}{rcl}
    \tilde S & \longrightarrow & \Sigma \\
    R_{\tilde s} & \longmapsto & \mathbb{C}\cdot R_{\tilde s} \in \mathbb{P}(V_{\mathbb{C}})(\ldots) = \mathbb{P}(V)(\mathbb{C}) .
  \end{array}
\]
\struck{Directement} \struck{\ill{}} En fait, \add{la construction $\Sigma$ provient} de la droite projective \emph{réelle} $\mathbb{P}(\underbrace{V_{\tilde P_{2n}}}_{V})$, et l'application (3) se factorise par \struck{\ill{}}
\[
  \text{(4)}\qquad \Sigma_{\mathbb{R}} = \Sigma_{\tilde P_{2n}}(\mathbb{R})
\]
où on a posé
\[
  \text{(5)}\qquad \underset{\textstyle \Sigma_{P_n}}{\Sigma_{\tilde P_{2n}}} = \mathbb{P}(V_{\tilde P_{2n}})
  \qquad \text{(ne dépend en fait que de } P_n = \tilde P_{2n}/\pm 1\text{)}
\]
\note{sous $\Sigma_{\tilde P_{2n}}$ dans (5), un signe d'égalité vertical renvoie à $\Sigma_{P_n}$.}

\page{96}
La représentation naturelle de \struck{$D_{\tilde P}$}
\[
  \text{(6)}\qquad D_{\tilde P_{2n}} \overset{\text{déf}}{=} \operatorname{Aut}_{\text{comb}}(\tilde P_{2n}) \qquad (\simeq \mathbb{D}_{2n})
\]
sur $V_{\tilde P_{2n}}$ définit une action de $D_{\tilde P_{2n}}$ sur $\Sigma_{\tilde P_{2n}}$ et par là sur $\Sigma$, mais qui est triviale sur le sous-groupe $\pm 1$, et se factorise donc par une action de
\[
  \text{(7)}\qquad D_{P_n} \overset{\text{déf}}{=} \operatorname{Aut}_{\text{comb}}(P_n) \simeq D_{\tilde P_{2n}}/\{\pm 1\}
\]
sur $\Sigma = \Sigma_{P_n}$, qui n'est autre que son action naturelle par transport de structure, quand on interprète $\Sigma$ directement comme défini en termes de $P_n$.

Si $n$ est impair, la suite exacte canonique
\[
  \text{(8)}\qquad 1 \longrightarrow \{\pm 1\} \longrightarrow D_{\tilde P_{2n}} \longrightarrow D_{P_n} \longrightarrow 1
\]
est scindée \struck{\ill{}} (le scindage \add{choix d'un} dépendant du choix d'un « sous-$n$-gône » \add{des deux} $P'_n$, $P''_n$ du $2n$-gône combinatoire $\tilde P_{2n}$ (pour $n = 3$, le choix d'un des deux triangles $ABC$, $A'B'C'$ marqués sur la figure). Cela est lié au fait qu'on a une équivalence de catégories, \emph{pour $n$ impair}
\drawing{en marge, un hexagone régulier inscrit dans un cercle, de sommets $A$, $C'$, $B$, $A'$, $C$, $B'$ ; les deux triangles $ABC$ et $A'B'C'$ y sont tracés.}

\page{97}
\note{p. 25 de l'auteur.}
\[
  \text{(9)}\qquad
  \begin{array}{ccc}
    \begin{array}{c}\text{Polygones combinatoires $\tilde P_{2n}$}\\ \text{d'ordre $2n$, avec un}\\ \text{sous-polygone d'ordre}\\ \text{$n$ marqué, soit $P'_n$}\end{array}
    & \xrightarrow{\ \approx\ } &
    \begin{array}{c}\text{polygones}\\ \text{combinatoires $P_n$}\\ \text{d'ordre $n$}\end{array} \\[2pt]
    (\tilde P_{2n}, P'_n) & \longmapsto & P'_n
  \end{array}
  \qquad (n \text{ \emph{impair}})
\]
\struck{Donc} l'\struck{représentation} \add{opération} naturelle de $D_{P_n}$ sur $\Sigma$, ou mieux sur $\Sigma_{\tilde P_{2n}}$ \add{$= \mathbb{P}(V_{\tilde P_{2n}})$}, se relève en une \struck{représentation} \add{opération linéaire} de $D_{P_n}$ sur $V_{\tilde P_{2n}}$. En fait, \struck{on peut même prendre donc} \add{on a}
\[
  \text{(10)}\qquad V_{\tilde P_{2n}} = \underbrace{V_{P'_n} \simeq V_{P_n}}_{\text{car } P'_n \simeq P_n}
\]
et l'action de $D_n$ sur $V_{\tilde P_{2n}}$ peut s'identifier à son action sur $V_{P_n}$. \struck{\ill{} $\Sigma_{\tilde P_{2n}}$ Ici $\Sigma_{\tilde P_{2n}}$ s'identifie à la droite projective} L'isomorphisme (10) fournit d'ailleurs
\[
  \text{(11)}\qquad \underset{\textstyle \Sigma_{P_n}}{\Sigma_{\tilde P_{2n}}} = \mathbb{P}(V_{\tilde P_{2n}}) \xrightarrow{\ \sim\ } \mathbb{P}(V_{P_n}) , \qquad (n \text{ \emph{impair}})
\]
\struck{et ce qui} et cette explicitation de $\Sigma_{P_n}$ comme $\mathbb{P}(V_{P_n})$ est \struck{\ill{}} d'ailleurs évidente a priori, car l'application
\[
  \begin{cases}
    S = S_{P_n} \longrightarrow \mathbb{P}(V_{P_n})(\mathbb{R}) \\
    s \longmapsto \mathbb{R}s
  \end{cases}
\]
est \emph{injective} pour $n$ impair — de sorte que

\page{98}
$\mathbb{P}(V_{P_n})(\mathbb{C})$ satisfait \struck{\ill{}} la propriété caractéristique de $\Sigma_{P_n}$, comme \uncertain{enveloppe}-droite-projective de \struck{\ill{}} $S_{P_n}$, « stable » compatible avec l'action de $D_{P_n}$.

Si par contre $n$ est pair, l'homomorphisme \struck{\ill{}} $D_{\tilde P_{2n}} \to D_{P_n}$ ne se scinde pas, car l'homomorphisme
\[
  \begin{array}{ccc}
    D_{\tilde P_{2n}}^+ & \longrightarrow & D_{P_n}^+ \\
    \wr & & \wr \\
    \mathbb{Z}/2n & & \mathbb{Z}/n
  \end{array}
\]
ne se scinde pas (il est isomorphe à l'homom. de passage au quotient canonique $\mathbb{Z}/2m \to \mathbb{Z}/m$, qui se prolonge au quotient par \struck{\ill{}} $m\cdot\mathbb{Z}$ (où $n = 2m$) dans $\mathbb{Z}/4\mathbb{Z} \to \mathbb{Z}/2\mathbb{Z}$, qui n'est pas scindée. \struck{Donc} Comme l'extension (8) de $D_{P_n}$ par $\{\pm 1\}$ n'est autre que l'extension réciproque de
\[
  1 \to \{\pm 1\} \to \operatorname{Sl}(V_{\tilde P_{2n}}) \to \underset{\textstyle \operatorname{Aut}(\Sigma_{P_n})}{\operatorname{GP}(V_{\tilde P_{2n}})} \longrightarrow 1
\]
par $D_{P_n} \to \operatorname{Aut}(\Sigma_{P_n})$, cela montre que l'opération de $D_{P_n}$ sur $\Sigma_{P_n}$ ne se remonte pas en une opération vectorielle — \struck{\ill{}} \add{\ill{}} qui reviendrait au même, qu'elle ne se remonte pas

\page{99}
\note{p. 26 de l'auteur.}
en une opération de $D_{P_n}$ sur $(\Sigma_{P_n}, \mathcal{O}_{\Sigma_{P_n}}(1))$ où $\mathcal{O}_{\Sigma_{P_n}}(1)$ désigne un faisceau inversible sur $\Sigma_{P_n}$, de degré 1 (p.ex. celui provenant de la description \struck{\ill{}} $\Sigma_{P_n} \simeq \mathbb{P}(V_{\tilde P_{2n}})(\mathbb{C})$). Si donc on s'intéresse au fibré principal
\[
  \text{(12)}\qquad P\Sigma_{P_n} = \check{\mathbb{V}}(\mathcal{O}_{\Sigma_{P_n}}(1))^*
\]
sur $\Sigma_{P_n}$, il n'y a pas d'action naturelle de $D_{P_n}$ sur ce fibré — mais (une fois choisi le revêtement $\tilde P_{2n}$ de degré 2 de $P_n$, ce qui précise le choix de $\mathcal{O}_{\Sigma_{P_n}}(1)$ en termes de $V_{\tilde P_{2n}}$), on trouve une action de $D_{\tilde P_{2n}}$ sur (12).

Dans tous les cas, $n$ impair ou pair, je m'intéresse : l'action de $D_{\tilde P_{2n}}$ sur (12), et sur \struck{le groupe}
\[
  \text{(13)}\qquad P\Sigma_{P_n}^* = P\Sigma_{P_n} \mid \Sigma_{P_n}^* ,
\]
qui est un fibré principal de groupe $\mathbb{C}^*$ sur $\Sigma_{P_n}^*$. De façon précise, je m'intéresse au groupoïde fondamental de $P\Sigma_{P_n}^*$ et l'action

\page{100}
de $D_{\tilde P_{2n}}$ dessus — et en particulier, au groupe fondamental mixte de $P\Sigma_{P_n}^*$, pour cette action de $D_{\tilde P_{2n}}$ \struck{\ill{}}, ou \ill{} pour l'action de $D_{\tilde P_{2n}} \times \{1,\tau\}$. \struck{\ill{}}

Quand
\note{la page s'arrête sur ce mot ; la suite manque dans le lot.}
\end{document}
